% Quantitative refinement of the shared source and compression proofs.
\section{Polynomial sufficient widths for Legendre and Taylor}
\label{cp:poly:width-section}

Use the common setup of \Cref{cp:setup} and, locally, put
\[
X=\beta^L,\qquad \lambda=\gamma/m\in(0,1],\qquad
\chi=1+\lambda^{-1},\qquad S=16Y/\lambda,\qquad \ell_n=\log(en).
\]
Thus $X\ge100$, $Y/\lambda\le X^{-30}$ and $S\le X^{-26}$;
$S$ is a normalization, not state. For $0<\rho<1/4$ and $N\ge m$
declared inputs define
\[
u_\rho=\max\left\{1,\left\lceil
 \frac{\log(4emL/\rho)}{\log(64e^2)}\right\rceil\right\},\qquad
A_{N,\rho}=X^{4000}\chi^8(N+d+u_\rho+1)^8,
\]
\begin{equation}
N_{\rm src}(N,\rho)
=\left\lceil(A_{N,\rho}/\rho)^{40000}\right\rceil.
\label{cp:poly:eq:source-width}
\end{equation}
This deliberately loose finite gate contains no $Y^{-1}$.

\begin{theorem}[Polynomial onset with the original storage powers]
\label{cp:poly:width}

For $0<\delta,\varepsilon<1$, a sufficient width for

\[
\Pr\left\{
 \|f_{\rm model}-f_n\|_{\mathcal X}
 \le\varepsilon\|f_n-\widetilde f_n\|_{\mathcal X}
 \right\}\ge1-\delta
\]

is bounded by

\[
\left[c_{\rm wid}\beta^L
 (1+m+d+p+\gamma^{-1}+Y^{-1}+\delta^{-1}+\varepsilon^{-1})
 \right]^{c_{\rm wid}}
\]

with one absolute finite numerical constant $c_{\rm wid}$, local to this theorem. For Legendre take $p=0$ and omit
$Y^{-1}$; its domain is the whole input sphere. For Taylor, $p$ is the
number of passive inputs declared before initialization and its domain is
that panel together with the training inputs. The Legendre and Taylor
storage bounds in \eqref{cp:legendre-size} and
\eqref{cp:panel-size} hold with activation/depth coefficients bounded
by fixed powers of $X$. Taylor also has absolute error $Y/n$.

The constructions, physical optimizer, initialization-only provenance and
label condition are unchanged. Confidence applies to each individual width.
The inverse-label gate for Taylor is needed only to simplify its exact
source degree to a multiple of $\log(en)$. This result makes no polynomial
onset claim for Harmonic, for the nonlinear feature-learning certificate,
or for the stronger coefficient in \eqref{cp:eq-dense-lower}, and gives no
pointwise-in-time or endpoint denominator bound.

\end{theorem}

The proof quantifies the shared source construction, derives a finite dense
lower bound, and improves the Legendre history defect to order $q^{-3}$.

\subsection{Finite initialization and analytic source probability}
\label{cp:poly:source-section}

Use exactly the local coefficient recurrences
\cref{cp:src-basic,cp:src-absorption-coefficients,cp:src-allowance,
cp:src-response-coefficients,cp:src-UV,cp:src-K}, with
\[
b=\max_j|\phi_j(0)|,\quad
s=\max\{1,\sup_{j,|\operatorname{Im}z|\le a/2}|\phi'_j(z)|\},\quad
t_2=\max\{1,\sup_{j,|\operatorname{Im}z|\le a/2}|\phi''_j(z)|\}.
\]
In particular $C_{\rm abs},W_G$ keep
\eqref{cp:src-absorption-coefficients}'s definitions.
The carriers are $k_a^{(L)}=w$,
$k_a^{(j)}=W^{(j+1)\top}\delta_a^{(j+1)}$ for $j<L$.
Write $\rho(t)=\|r(t)\|_2/\sqrt m$ for residual RMS, distinguishing it
from the constant failure allowance $\rho$, and use Euclidean mobility
coordinates
$\theta=(W^{(1)}/\sqrt n,W^{(2)},\ldots,W^{(L)},w/\sqrt n)$.
Only finite input lists are used in the complex source argument.

\begin{lemma}[Finite training and panel sources]
\label{cp:poly:source}

At $n\ge N_{\rm src}(N,\rho)$, with probability at least $1-\rho$,
the ambient real fitting event holds and the following source conclusions
hold. The four families are exactly \eqref{cp:source-families};
forward inputs range over the $N$ declared inputs and backward inputs over
the $m$ training inputs.

For the training list alone one may use the closed rectangle

\begin{equation}
[-r,T+r]+i[-r,r],\qquad
T=32\ell_n/\lambda,\qquad
r=[u_\rho X^{100}(1+\lambda)\sqrt{\ell_n}]^{-1}.
\label{cp:poly:eq:training-radius}
\end{equation}

All training sources and the responses
$R_{ba}^{(j)}=D_\Theta z_b^{(j)}\nabla_\Theta(nf_a)$, where

\[
\Theta=(W^{(1)},\sqrt nW^{(2)},\ldots,\sqrt nW^{(L)},w),
\]

are holomorphic near this rectangle. The source RMS bounds are
\eqref{cp:source-rms}, their coordinates are bounded by
$X^6\sqrt n$, and $|f_n|\le16X^6Y/\lambda$. Moreover

\[
\max_{a,b,j,i}|R_{ba,i}^{(j)}|\le SU_{\rm fin}\sqrt{\ell_n}
\]

on that rectangle. On either complex source domain the first-layer
operator divided by $\sqrt n$ and all hidden operators have norm at most
ten, the complex training residual RMS is at most $2Y$, and every declared
preactivation has imaginary part below $3a/8$. The training coordinate
maxima are bounded by $K_{\rm src}\sqrt{\ell_n}$ for preactivations and
$SK_{\rm src}\sqrt{\ell_n}$ for carriers. These strict strip and
coordinate bounds are included in the event, not inferred from an RMS
bound. On the whole real training trajectory, including its
limit,

\begin{equation}
\max_{a,j,i}|k_{a,i}^{(j)}|\le32X^{21}(Y/\lambda)\sqrt{\ell_n},
\qquad
\max_{a,b,j,i}|R_{ba,i}^{(j)}|
 \le YX^{110}\chi^3(1+\sqrt{\ell_n}).
\label{cp:poly:eq:real-response}
\end{equation}

For the $N$-input panel put $a_0=\min(a,1)$ and

\begin{equation}
b_n=\log(e+\ell_n),\quad B_n=1+(S^2/\eta_*)b_n,\quad
r_n=\min\left\{\frac{a_0}{64YSU_{\rm fin}\sqrt{\ell_n}},
 \frac1{X^{300}(1+\lambda B_n)\sqrt{b_n}}\right\}.
\label{cp:poly:eq:panel-radius}
\end{equation}

There is a deterministic disk cover and half-radius panel partition of
$[0,T]$, with the same source envelopes, whose number of panels obeys

\begin{equation}
J\le1+X^{400}\left[(Y/\lambda)^2\sqrt{\ell_n}
                      +\lambda^{-1}\log(e+\ell_n)\right].
\label{cp:poly:eq:panel-count}
\end{equation}

Neither this radius nor its panel-count coefficient depends on confidence.

\end{lemma}

\begin{proof}

First quantify initialization. Let
$H=\max_{0\le j\le L}\sqrt{Q^{(j)}_{11}}\ge1$ and let $b,s$
have their meanings in the source proof. Define

\[
R_0=2H,\quad K_0=b+sR_0,\quad D_0^{\rm cov}=2sb+4s^2R_0,
\quad T_1=\sum_{i=0}^{L-1}s^i,\quad T_2=\sum_{i=0}^{L-1}s^{2i},
\]

\[
C_0^{\rm cov}=T_2(2K_0+1+D_0^{\rm cov}T_1),\quad
e_0=\min\{(8T_1)^{-1},\gamma/(4mC_0^{\rm cov})\},\quad
h_0=\min\{1/2,H/[4(8s)^L]\}.
\]

The Gaussian Lipschitz inequality in \Cref{cp:src-Gaussian} gives
$\operatorname{Var}F\le a^2$ and
$0\le(\mathbb EF^2)^{1/2}-\mathbb EF\le a$ for nonnegative
$a$-Lipschitz $F$, hence
\[
\Pr\{|F-(\mathbb EF^2)^{1/2}|>e\}\le2e^{-e^2/(8a^2)}
\qquad(e\ge2a).
\]
For iid Gaussian row pairs of marginal deviations at most $R_0$,
the empirical RMS of $\phi(Z_a)$ and $\phi(Z_a)\pm\phi(Z_b)$
has Lipschitz constant $sR_0/\sqrt n$ and $2sR_0/\sqrt n$,
respectively, also for singular pairs. Polarization therefore gives
entry error $(2K_0+1)e_0$ when the two latter RMS errors are at most
$e_0\le1$.

To propagate covariance errors without dividing by a variance, consider
two Gaussian pairs with standard deviations $\sigma_a,\sigma_b$ and
$\sigma'_a,\sigma'_b$, all at most $R_0$, and set
$\Delta=\max_i|\sigma_i-\sigma'_i|$.
Shrink both pairs to marginal deviations
$\tau_i=\min(\sigma_i,\sigma'_i)$, preserving their correlations.
Product subtraction and Cauchy--Schwarz bound the total scaling cost by
$2sK_0\Delta$. At these common marginals Gaussian integration by parts
gives the correlation derivative
$\tau_a\tau_b\mathbb E[\phi'(Z_a)\phi'(Z_b)]$.
The formula follows by differentiating the nonsingular density and
integrating its mixed second derivative; linear growth and bounded
derivatives justify the integrations. Continuity covers correlations
$\pm1$ and zero marginal variances. Since

\[
\tau_a\tau_b|\rho_c-\rho'_c|
 \le|C_{ab}-C'_{ab}|+2R_0\Delta,
\]

the transformed entry error is at most
$s^2|C_{ab}-C'_{ab}|+D_0^{\rm cov}\Delta$.
To check the inequality, put the correlation difference over the larger
of $\sigma_a\sigma_b,\sigma'_a\sigma'_b$ and use
$|C'_{ab}|\le\sigma'_a\sigma'_b$.

Condition at each layer on the preceding initialized layers. On a sphere
$h_0$-net of size at most $(1+2/h_0)^d$, and the training list, let
$\Delta_j$ denote the largest RMS error and $E_j$ the largest training
Gram-entry error. The preceding estimates give

\[
\Delta_j\le e_0+s\Delta_{j-1},\qquad
E_j\le(2K_0+1)e_0+s^2E_{j-1}+D_0^{\rm cov}\Delta_{j-1}.
\]

Thus $\Delta_j\le e_0T_1\le1/8$ and
$E_j\le e_0C_0^{\rm cov}\le\gamma/(4m)$. This also keeps all
conditional standard deviations below $R_0$.
With

\[
\Xi=\log\frac{16L(m+1)^2}{\rho_0}
                 +d\log(1+2/h_0),
\]

a sufficient initialized fitting width at failure $\rho_0$ is

\begin{equation}
\left\lceil\max\left\{1,
 \frac{\log(8L/\rho_0)}{8-2\log9},
 \frac{d\log9+\log(8/\rho_0)}{8-\log9},
 \frac{32s^2R_0^2\Xi}{e_0^2}\right\}\right\rceil.
\label{cp:poly:eq:initialization}
\end{equation}

The at most $3m^2+(1+2/h_0)^d$ norm tests per layer cost
$\rho_0/2$; the two operator-net tails proved in \Cref{cp:fit}
cost the other half. On that operator event the layer-$j$ feature map
has RMS Lipschitz constant $(8s)^j$, so off-net increments are at
most $H/4$. Thus \eqref{cp:eq-initialization-event} holds and its
deterministic fitting proof applies. Unrolling gives
\[
R_0,K_0\le X^3,\quad D_0^{\rm cov}\le X^4,\quad
T_1\le X^2,\ T_2\le X^3,\ C_0^{\rm cov}\le X^{10},\quad
e_0^{-1}\le X^{11}\chi.
\]
Since $32s^2R_0^2\le X^8$ and
$\log(1+2/h_0)\le3L\log\beta$, the width is at most
$\lceil X^{32}\chi^2[dL\log\beta+\log(16L(m+1)^2/\rho_0)]\rceil$,
without ordering $m,d$.

For source initialization take $\rho_0=\rho/64$.
Conditional Gaussian tails give
$Z_0=2H\sqrt{2\log(64nNL/\rho)}$ on the declared list.
Zero-embedding and forward subtraction give
$D_{\rm init}=(8s)^L\sqrt u(b+sZ_0)$ for deletion count
$u=u_\rho$. Since the full training Gram starts above $3\gamma/4$
and $D_{\rm init}/\sqrt n<\sqrt\lambda/8$, every cavity's normalized
least singular value exceeds
$(\sqrt3/2-1/8)\sqrt\lambda>\sqrt{\lambda/2}$.
Project lower differences before conditioning on the fresh row:
their $n^{-1/10}$ tail is
$2\exp[-n^{4/5}/(2D_{\rm init}^2)]$, with at most $uLn^uNLn$ tests.
Initial budget means and second moments are below four and fifteen,
since $\mathbb E e^{\eta_*|Z|}\le2e^{1/2}$,
$\mathbb E e^{2\eta_*|Z|}\le2e^2$ and $\eta_*2H\le1$.
Conditional Chebyshev costs $15mL/(16n)$ for layer means below eight.
All initialization failures total less than $\rho/4$ at the source gate.
Each cavity has its own initial test and zero reference paths on failure.

We now quantify the independently stopped source proof
\cref{cp:src-local,cp:src-local-lemma}. Use
$\bar u=(\Theta-\Theta_0)/S$, $\bar F_a=nf_a/S$ in
\eqref{cp:src-retained}. Since $w(0)=0$,
$\bar F_a=\bar u_w^\top h_a^{(L)}$ has no inverse-$S$ derivative.
Freeze the scalar forward and scaled reverse controls; concatenate omitted
interior roots into $g\sim N(0,I_r/n)$, $r\le2un$.
First-layer incoming rows enter only their own scalar controls, with no
retained reverse port.

Put $\kappa_0=1/4000$, $\Gamma=\sqrt{\mathcal K}$,
$\tau_*=\max_j\tau_j$, $f_*=\max_j f_j$, and $T_0=64/\lambda$.
The propagator bound \eqref{cp:src-propagator} is
$\|J(t,s)\|\le J_0n^{\kappa_0}$ with
$J_0=2e^{1/4}(2\mathcal B)^{1/4000}$.
Define
\[
M_0=\eta_*^{-1}[1+\log(2\mathcal B)],\quad
A_0=b+sM_0,\quad B_0=sM_0,\quad E_0=A_*+SD_*M_0.
\]
Budgeted scaled carriers are at most $M_0\ell_n$ and controls at most
$A_0\ell_n,B_0\ell_n$.
The exact product recursions in \Cref{cp:src-local-lemma}
give the following coefficient ledger; each column omits the common
factor $\sqrt u n^{\kappa_0}$:
\[
\begin{array}{c|ccc}
 &\text{parameter}&\text{forward}&\text{backward}\\ \hline
\text{map}&X^{80}\chi\ell_n^2&X^{85}\chi\ell_n^2&
                              X^{125}\chi\ell_n^3\\
\text{time derivative}&X^{125}\chi\sqrt n\ell_n^4&
 X^{140}\chi\sqrt n\ell_n^4&X^{190}\chi\sqrt n\ell_n^4
\end{array}
\]
Control differences have coefficient $X^{90}\chi\ell_n^2$ times
the supremum control difference. Here is a direct verification,
so none of these finite powers comes from a qualitative threshold.
The source recurrences and geometric sums give
\[
\begin{gathered}
H_j,P_j,f_j\le X^3,\quad k_j,\tau_j\le X^4,\quad V_j^0\le X^7,\\
A_*,D_*,\eta_*^{-1},\mathcal B,\Gamma,J_0
 \le X^5,X^6,X^{27},X^3,X^7,X,\\
M_0,A_0,B_0,E_0\le X^{31},X^{33},X^{32},X^{38}.
\end{gathered}
\]
The parameter map coefficient supplied by the three port derivatives is
\begin{equation}
v_0=J_0\{A_0(2T_0\Gamma\tau_*+E_0)+Sf_*B_0\}
       \le X^{80}\chi.
\label{cp:poly:eq:map-coefficients}
\end{equation}
Duhamel uses $2\int\rho/S\le1$ and length $T_0\ell_n$.
Forward coefficients are $P_j(Sv_0+A_0)$, multiplied by $s$
for features. Backward coefficients start at $v_0$ and satisfy
$k_j^0=10d_{j+1}^0+S\tau_{j+1}v_0+B_0$,
$d_j^0=sk_j^0+t_2M_0P_j(Sv_0+A_0)$.
Their only carrier factor is additive. Using
$(10s)^L\le X^2$ and at most $L\le X$ forcing terms proves
the first row. Replacing $A_0,B_0$ by one in these forcings gives
the stated control bound, without changing reference coefficients.

For time moduli, differentiating the reference recursion gives
\[
T_L^0=2sH_L+2t_2M_0V_L^0,\qquad
T_j^0=10sT_{j+1}^0+2s\tau_{j+1}^2H_j+2t_2M_0V_j^0.
\]
Thus $T_j^0\le X^{45}$ and the control speed is at most
$D_c\sqrt n\ell_n$, where
$D_c=\max\{\lambda s\max_jV_j^0/4,\lambda\max_jT_j^0/8\}
\le\lambda X^{45}$.
All reference derivative coefficients, before the single raw-coordinate
factor $\sqrt n\ell_n$, are at most $\lambda X^{47}$.
The time equation for the parameter map has coefficient
\[
v_t=[2\mathcal K+(\lambda/4)SE_0]v_0
 +A_0(2\Gamma\tau_*+\lambda E_0/4)+(\lambda/4)Sf_*B_0
 \le X^{125}\chi.
\]
In the differentiated forward recursion, the largest new coefficient
is bounded by $X^{47}X^{85}\chi$; adding the gate derivative and
geometric propagation gives $X^{140}\chi$.
For backward derivatives the forcing
$t_2M_0X^{140}\chi$ is at most $X^{172}\chi$.
The other gate terms are $t_2\dot z\,k_{[1]}$ and
$(\phi'''\bar k\dot z+\phi''\dot{\bar k})z_{[1]}$.
With $|\phi'''|\le4\beta^2$ and the reference speed bound, they are
at most $X^{174}\chi$; propagation and all sums fit below
$X^{190}\chi$.
Differentiating the lower-Jacobian probes uses the same forward
products. Every term has at most one raw coordinate derivative,
hence one factor $\sqrt n$; the log powers are at most four.
This proves the second row without differentiating the variational Jacobian.
Consequently $C_0=X^{300}\chi$ bounds every coefficient above.
For a root selector $E_i$, $\operatorname{sym}(E_i^\top L)$
inherits the linear map's moduli; centered interpolation on
$\|g\|\le2\sqrt{2u}$ costs $10u\|\Delta L\|$, including trace change.

For responses use \eqref{cp:src-response-recursion}, with evaluated
input $b$ and training driver $a$, $\bar R=R/S$ and
$\bar Q=\phi'\odot\bar R$. Its linearized products and the direct
reverse observable $C_bC_a^\top\bar q_a$, $C_a=D_\Theta h_a$,
were proved in \Cref{cp:src-local-lemma}. In particular
\[
(C_b^1C_a^{1\top}-C_b^0C_a^{0\top})y_i
=(C_b^1-C_b^0)C_a^{0\top}y_i
 +C_b^1(C_a^1-C_a^0)^\top y_i
\]
retains $d_0+n^{-1/2}$ by the reference-probe and lower-adjoint bounds
there. For time derivatives the extra product is
$\phi''\dot z_b\bar R$: use the response coordinate cap and $z_b$'s
RMS speed before one scalar-control conversion. The raw amplitude
coefficient is at most $X^{80}$ and $U_{\rm fin}\le X^{26}$, so
the same product recursions have common envelope
\[
X^{100}(1+X^{80}+4U_{\rm fin})^2(1+C_0)^2
 \le X^{1000}\chi^2=:\mathcal C.
\]
The exponent count is $100+160+600$, with the remaining $140$ paying
all numerical sums and bounded layer propagation. Thus the map norms,
control moduli and actual analytic derivatives are bounded by

\begin{equation}
\mathcal C\sqrt u n^{\kappa_0}\ell_n^5,\qquad
\mathcal C\sqrt u n^{\kappa_0}\ell_n^4,\qquad
\mathcal C\sqrt{un}n^{\kappa_0}\ell_n^8.
\label{cp:poly:eq:map-interface}
\end{equation}

There are at most $4u(N+1)^2$ real scalar controls, of amplitude
$M=\mathcal C\ell_n$ and speed $D=\mathcal C\sqrt n\ell_n^3$.
Every bound concerns frozen one-dimensional histories until the uniform
event is established.

We check the deterministic domain premises before the Gaussian event.
The training radius is at most
\[
\min\{1/8,1/\lambda,\log2/(4\mathcal K),1/(32YSD_W),
                a/(64YSU_{\rm fin}\sqrt{\ell_n})\}.
\]
Thus the residual, activity, operator and pole estimates in
\hyperref[cp:src-whole-sphere-verification]{the rectangle verification}
apply without angular fields or an $n^{-1}\le Y$ gate.
For the panel use exactly $h,d_h,T_+,D_s,P_s$ of
\Cref{cp:src-panel-working}, replacing its radius by
\eqref{cp:poly:eq:panel-radius} and taking $\nu=\lambda/4$.
That construction gives $\operatorname{Lip}P_s\le1+d_h$,
$d_h\le1/8$, commuting real anchors, a bounding rectangle of sides
at most $H_0=128\ell_n/\lambda$, and canonical contours of length
at most $H_0/2$. The recurrences
\cref{cp:src-time-derivative,cp:src-gradient-derivative} give
$J_*\le X^{34}$, $C_g\le X^{38}$, hence
\[
4\mathcal Kr_n+(2C_g/\sqrt{\mathcal K})B_nYr_n
 \le4\mathcal K/X^{300}+C_g/(8X^{300})<1/100,
\]
using $Y=\lambda S/16$ and $\lambda B_n/(1+\lambda B_n)\le1$.
This verifies the complete propagator gate
\eqref{cp:src-flaring-gate}: the frozen real generator is unitary
in imaginary time, its perturbation costs the second term, and forward
real segments contract. The same domain proof gives additional activity
$8Yr_n$, normalized operator motion $8YSD_Wr_n$, and improved
vertical preactivation displacement
$2YSU_{\rm fin}\sqrt{\ell_n}r_n\le a_0/32$.
The negative corner costs at most $a_0/8$; all margins are strict.

The uniform event must precede adaptive substitution; endpoints are
clamped only by their cavity's own stop. Put
$A=A_{N,\rho}=\mathcal C^4(N+d+u+1)^8$. The source gate gives
\begin{equation}
A\ell_n^{16}\le n^{1/1000},\quad
u\le n^{1/8000},\quad \ell_n\le n^{1/16000}.
\label{cp:poly:eq:absorption-gates}
\end{equation}
Indeed, with $v=\log(A/\rho)\ge8000\log10$ and
$\log n=40000v$, one has
$\log A+16\log(1+40000v)\le17v<40v$; the difference increases
thereafter. The other two inequalities use the same gate.

Condition on a cavity's retained initialization, with own stop $\sigma_c$.
Set $\xi=n^{-2}/(1+A\ell_n^8)$,
$J_s=\lceil16H_0/\xi\rceil$, $s_j=j/J_s$, and let $\mathscr G$
be a product grid of spacing at most $\xi/8$ on the bounding rectangle.
Test the retained-measurable endpoints
\begin{equation}
z_{j,q,c}=P_{\min(s_j,\sigma_c)}q,
\qquad q\in\mathscr G,\quad0\le j\le J_s.
\label{cp:poly:eq:prefix-grid}
\end{equation}
For endpoint $x+iy$ its own-domain contour is
$\zeta(t)=\operatorname{sgn}(x)\min(t,|x|)
+i\operatorname{sgn}(y)\min((t-|x|)_+,|y|)$ on $[0,H_0]$,
then frozen. Solve
$V'=\zeta'(\mathsf AV+\sum_a\mathsf B_a^ee_a+\mathsf B_a^q\bar q_a)$,
$V(0)=0$, with the preceding port derivatives. Bounded measurable
histories suffice, and $V$ is a linear map of omitted roots; its
clamped extension is not claimed holomorphic.

For control accuracy $\tau=n^{-1/8}$ sample $[0,H_0]$ in
$\lceil16DH_0/\tau\rceil$ equal intervals, round each real component
to mesh $\tau/8$ and interpolate. Every admissible history is within
$\tau$ and the logarithmic cardinality is at most

\[
4u(N+1)^2(2+16DH_0/\tau)\log(1+32M/\tau)
 \le A\ell_n^8n^{5/8}\le n^{0.627}.
\]

Here $16384u(N+1)^2\mathcal C\lambda^{-1}\le A$.
The net amplitude is at most $2\mathcal C\ell_n$; affine control
dependence and \eqref{cp:poly:eq:map-interface} give map norm
$4\mathcal C^2\sqrt u n^{\kappa_0}\ell_n^5
\le n^{1/800}<n^{1/200}$.

For $g\sim N(0,I_r/n)$, $r\le2un$, a real symmetric matrix obeys

\[
\Pr\{|g^\top Rg-\operatorname{tr}R/n|>t\}
 \le2\exp\left[-\frac18\min\left\{
 \frac{n^2t^2}{\|R\|_{\rm HS}^2},\frac{nt}{\|R\|}\right\}\right].
\]

Diagonalization and
$-\log(1-v)-v\le v^2/[2(1-|v|)]$, followed by exponential Markov
at parameter at most $n/(4\|R\|)$, prove this inequality as in
\Cref{cp:src-local-lemma}. Since
$\|R\|_{\rm HS}^2\le2un\,n^{1/100}$, real/imaginary splitting
at $n^{-1/10}/8$ costs at most $4e^{-n^{0.78}}$.
Linear-coordinate tails are smaller. The Euclidean image mean is
at most $\sqrt{2u}n^{1/200}<n^{1/100}/4$; Gaussian concentration
pays the remaining threshold $n^{1/100}/2$.
Root-norm failure is at most $e^{-un}$, and control interpolation
costs $10\mathcal C u^{3/2}n^{\kappa_0-1/8}\ell_n^4<n^{-1/10}/8$.

For an arbitrary full-network-dependent common prefix $D_\sigma$,
choose a nearest $q\in\mathscr G$ and the largest $s_j\le\sigma$.
Since $\sup_z|P_sz-P_tz|\le3H_0|s-t|$,
\[
|z-P_{s_j}q|\le(3/16+9\sqrt2/64)\xi\qquad(z\in D_\sigma).
\]
Its canonical contour lies in $D_\sigma$, so actual controls may now
be substituted. The domain is vertically fibered and retracts to its
real interval, hence is simply connected. Local analytic ODE uniqueness
and continuation along compact paths give a single holomorphic variation
for these actual controls. The $P_\sigma$-image of the joining segment
has length below $\xi/2$; the derivative bound
\eqref{cp:poly:eq:map-interface} therefore costs
$10u\mathcal C\sqrt u n^{\kappa_0+1/2}\ell_n^8\xi<n^{-1/10}/8$.
Continuity covers boundary levels, and level zero is tested.
There are at most $(2+16H_0/\xi)^3\le n^9$ endpoints,
$An^3$ maps per endpoint and $uLn^u$ deletion sets. Their logarithmic
union cost is below $n^{0.65}$, so failure is below
$e^{-n^{0.7}}<\rho/16$.

On rectangle prefixes, round each coordinate toward zero to mesh $\xi$:
the grid point and joining segment remain in every common prefix.
Clamp each canonical contour only by its cavity's own rectangle.
The same conditional tests and interpolation then require at most
$n^6$ endpoints and the same failure allowance.

Apply the exact remainder formulas
\cref{cp:src-exact-remainders,cp:src-remainder-bounds,cp:src-force-remainder}
with stops $N_0=n^{1/100}$, $d_0=n^{-1/10}$, $u_0=n^{-1/25}$ and
$R=d_0N_0+d_0u_0+u_0^2$, $P=(N_0+u_0)^2/\sqrt n$.
The force coefficient is $X^{240}(1+u)(1+M)$; learned ports and the
top offset cost $X^{260}(1+u)^2(1+M)^3$.
Their raw envelope $M\le X^{80}\ell_n$, the map ledger,
$2\int\rho/S\le1$ and length $64\ell_n/\lambda$ give
\[
\|U\|\le A\ell_n^8n^{\kappa_0}
 [n^{-0.08}+u_0^2+n^{-0.48}+n^{-1/2}]<u_0/2.
\]
The largest power is $n^{-0.07875}$, and the source gate also gives
$d_0+X^{10}u_0+X^{30}(R+P)<a/32$.
Thus the same first-exit proof gives coordinate differences
$n^{-1/30}$ and vector differences $2n^{1/100}$, including responses.
Its first-layer incoming increment $sS^2M$ belongs to deleted scalar
controls, with no retained reverse port.

The four traces \eqref{cp:src-traces}, the Schatten calculation and
scalar feedback solution in \hyperref[cp:src-moments]{the source proof}
apply before adaptive substitution. The first three constants are
unchanged, the residual trace is at most
$A\ell_n^8n^{-1+\kappa_0}$, and finite remainders are below $n^{-1/40}$,
so the original $C_{\rm abs}$ applies.
The Gaussian maximum argument there gives
$K_{\rm src}\sqrt{\ell_n}$ and $SK_{\rm src}\sqrt{\ell_n}$.
Its four response terms have, after division by $S$, precisely
\[
sK_{\rm src}H_{j-1}^2,\quad sK_{\rm src}f_{j-1}^2,\quad
sK_{\rm src}ST_Q,\quad sK_{\rm src}S^2H_{j-1}q_{j-1}^*.
\]
Adding the centered coefficient $64q_{j-1}^*$ gives $U_j(64)$.
Multiplier 64 pays the $mNLn$ response tests and polynomial endpoints
at the source gate; both maximum failures total less than $\rho/16$.
The exact identity $\dot z_b=-(2/m)\sum_a r_aR_{ba}$ improves poles
and passive maxima, with real displacement
$2S^2U_{\rm fin}\sqrt{\ell_n}$ and no passive budget.

Budget removal must hold at the confidence order $u=u_\rho$.
Complete each common-cavity reference by its own stopped retraction
before integrating a root. For normalized complex-minus-real Gaussian
coefficients the panel derivative estimate gives
\[
D\le X^{100}\lambda r_nB_n,\quad K\le X^{100}\lambda B_n,\quad
H\le128\ell_n/\lambda.
\]
Use $\bar D=X^{-200}\lambda B_n/[(1+\lambda B_n)\sqrt{b_n}]$.
By \Cref{cp:src-Gaussian}, the logarithmic moment is at most
$a(M_G+2\bar D)+2a^2\bar D^2$, with
$M_G=256\bar D\sqrt{\log(e+HK/\bar D)}$, and
\[
\log(e+HK/\bar D)\le300\log X+16b_n+\log(1/\lambda).
\]
For $b\ge1,x\ge0$,
$\min(1,be^{-x})\sqrt x\le\sqrt{\log b}+1$:
split at $x=\log b$ and use $e^{-s}\sqrt s\le1$ for the excess.
Taking $b=B_n$, $x=\log(1/\lambda)$ and $B_n\le1+b_n$ gives
$M_G+2\bar D\le X^{-190}$ and moment at most $e^{X^{-180}}$
for $0\le a\le8$, including sample RMS.

On the training rectangle take
$D\le X^{40}\lambda B_nr$, $K\le X^{40}\lambda B_n$,
$H\le64\ell_n/\lambda$. Its entropy is bounded by
$\log128+\log u+100\log X+\log\chi+(3/2)\log\ell_n$.
Substitution of \eqref{cp:poly:eq:training-radius},
$u^{-1}\sqrt{\log u}\le1$ and $\lambda\sqrt{\log\chi}\le2$
bounds the mean plus $2D$ by $X^{-50}$.
For the remaining logarithms use
$\log(e+x)\le2+\log x$ and maximize
$(\log x)^a/\sqrt x$, $a=0,1/2,1,3/2$:
this bounds
$(1+\log(e+\ell_n))(1+\sqrt{\log\ell_n})/\sqrt{\ell_n}$,
with constants absorbed by five of the sixty available powers of $X$.
Thus both geometries have small fixed-exponent corrections.

For each distinct root $i$, let $A_i$ exponentiate
$\eta_*C_{\rm abs}$ times one plus its four complete common-cavity
real scalar/sample-RMS references and four complex-minus-real corrections.
Equation \eqref{cp:src-main-moment} and Cauchy--Schwarz give
$\mathbb E(A_i^2\mid\text{cavity})<3$; Hölder among the four
corrections uses $16\eta_*C_{\rm abs}<8$.
The $A_i$ are conditionally independent because their root pairs are distinct.

Only singleton/common-cavity differences need order $u$.
Extend each path by its own stop and radially project their difference to
radius $4n^{1/100}$; both omit the paired root. Their Gaussian radius is
$v_n=4n^{-49/100}$ and \Cref{cp:src-Gaussian} gives mean
\[
\mu_n=256v_n\sqrt{\log(e+X^{400}\chi\ell_n^4/v_n)}+2v_n,
\qquad u\mu_n+u^2v_n^2<10^{-3}.
\]
Let $E$ be the successful initialization, map and maximum event above.
For $B_i^{\rm wt}=e^{\eta_*(Z_{a,i}+K_{a,i}/S)}$ and distinct
$|I|=k\le u$, scalar absorption, Cauchy--Schwarz between main references
and corrections, and Hölder over at most four corrections per root give
\[
\mathbb E\left[1_E\prod_{i\in I}B_i^{\rm wt}\right]
 \le3^{k/2}\exp\{kn^{-1/40}+4k\mu_n+64k^2v_n^2\}\le16^k.
\]
The indicator is dropped only after domination by complete references;
no sample-process or correction-path independence is used.

From the defining constants,
$2\eta_*K_{\rm src}\le33/4096<1/100$, so repeated weights cost at most
$W_n=e^{\sqrt{\ell_n}/100}$. Partition an ordered $u$-tuple by its
distinct indices. A partition with $j$ repetitions is encoded by $j$
edges from block minima, giving at most $\binom u2^j/j!$ possibilities.
The preceding distinct-product bound therefore yields

\[
\mathbb E\left[1_E(n^{-1}\sum_iB_i^{\rm wt})^u\right]
 \le16^u\exp\{\tbinom u2W_n/(16n)\}\le e16^u,
\]

where $n\ge4u^3$ suffices and is implied by the source gate.
Markov and the sample/layer union make the budget-hit probability at most
$emL(64e^2)^{-u}\le\rho/4$.
Cavity budgets transfer as
$\mathcal H_a^{\rm cavity}\le
e^{2\eta_*n^{-1/30}}\mathcal H_a+u/n<2\mathcal B$;
every other stop has a strict margin.
For either geometry, initialization, uniform maps, maxima and budgets
cost respectively $\rho/4,\rho/16,\rho/16,\rho/4$.
Initialization is shared, so doubling only the other three allowances
gives total failure below $\rho$. Holomorphic local existence and
uniqueness continue the solution across the compact domain boundary.

For real $t\ge T$, fitting gives
$\|\theta(t)-\theta(T)\|\le2Y\lambda^{-1/2}e^{-16\ell_n}$.
The backward subtraction of \Cref{cp:src-finite-working}, with coefficient
$X^{20}n$, gives the carrier cap. In response subtraction about $T$,
changed gates multiply reference coordinate caps, and changed mixers
and pairings use operator/RMS bounds. Its RMS Lipschitz coefficient is
$X^{100}\chi^2(1+\sqrt{\ell_n})$; one coordinate conversion leaves tail
$YX^{103}\chi^3(1+\sqrt{\ell_n})n^{-31/2}
\le Y\chi\sqrt{\ell_n}$ at the source gate.
This proves \eqref{cp:poly:eq:real-response}, including the limits.

The deterministic disk argument of \Cref{cp:src-panel-working}
uses $\varrho(t)=h(t)/[2(1+d_h)]$ and
$t_{b+1}=\min(T,t_b+\varrho(t_b)/2)$, giving
\[
J\le1+40/(\lambda r_n)
 +80\mathcal K\lambda^{-1}\log(1+4\ell_n/\log2).
\]
For the capped radius,
\[
\frac1{\lambda r_n}\le
\frac{1024U_{\rm fin}}{a_0}(Y/\lambda)^2\sqrt{\ell_n}
 +X^{300}\lambda^{-1}\sqrt{b_n}+X^{300}B_n\sqrt{b_n}.
\]
Writing $x=\log\ell_n\ge0$ gives $b_n\le2+x$ and
$(2+x)^{3/2}e^{-x/2}\le4$, so $b_n^{3/2}\le4\sqrt{\ell_n}$.
Together with $S^2=256(Y/\lambda)^2$, $a_0^{-1}\le\beta$
and the coefficient bounds this proves \eqref{cp:poly:eq:panel-count}.
\end{proof}

\subsection{A finite lower bound for actual dense variability}
\label{cp:poly:dense-lower-section}

For $0<\tau<1$, define

\[
a_\tau=\frac{\tau\gamma^2}
 {2^{20}m u_{\tau/8}X^{108}(1+\lambda)},
\]

\begin{equation}
N_{\rm low}(\tau)=\left\lceil\max\left\{\begin{gathered}
 N_{\rm src}(m,\tau/8),\quad
 2^{44}m^6X^{96}\gamma^{-12}\tau^{-2}\log(64/\tau),\\
 [2^{30}u_{\tau/8}X^{114}m\chi(1+\gamma^{-1})^2/\tau]^2
 \end{gathered}\right\}\right\rceil.
\label{cp:poly:eq:lower-width}
\end{equation}

\begin{proposition}[Quantitative trajectory separation]
\label{cp:poly:dense-lower}

For each $n\ge N_{\rm low}(\tau)$, two independent dense runs satisfy

\begin{equation}
\Pr\left\{\|f_n-\widetilde f_n\|_{\mathcal X}
       \ge a_\tau\frac{Y}{\sqrt n\ell_n^{5/2}}\right\}\ge1-\tau.
\label{cp:poly:eq:dense-lower}
\end{equation}

The domain may be any promised query set containing the training inputs.
The lower coefficient is weaker than that in \eqref{cp:eq-dense-lower}.

\end{proposition}

\begin{proof}

We first prove a shift-uniform finite scalar inequality. For iid centered
real variables $V_i$ with variance $\sigma^2>0$ and third absolute moment
at most $M_3$, put $Z=n^{-1/2}\sum_iV_i$. For every $h>0$ and real $b$,

\begin{equation}
\Pr\{|Z-b|\le h\}\le3h/\sigma
 +4\exp[-9h^2\sigma^4n/(128M_3^2)].
\label{cp:poly:eq:finite-smallball}
\end{equation}

For an independent copy $V'$, Taylor expansion gives
$|\mathbb Ee^{itV}|^2=\mathbb E\cos(t(V-V'))
\le1-\sigma^2t^2+(4/3)M_3|t|^3$, because
$\mathbb E|V-V'|^3\le8M_3$. Therefore
$|\mathbb Ee^{itV}|\le e^{-\sigma^2t^2/4}$ for
$|t|\le3\sigma^2/(8M_3)$.
The Gaussian majorant
$1_{|z-b|\le h}\le\sqrt e\,e^{-(z-b)^2/(2h^2)}$ and its Fourier
integral bound the probability by

\[
\frac{\sqrt e\,h}{\sqrt{2\pi}}
 \int_{\mathbb R}e^{-h^2t^2/2}|\mathbb Ee^{itZ}|\,dt.
\]

On $|t|\le3\sigma^2\sqrt n/(8M_3)$ use the preceding Gaussian
bound; its full-line integral is below $3h/\sigma$.
On the complement use modulus one and
$\Pr(|G|>s)\le2e^{-s^2/2}$. This proves
\eqref{cp:poly:eq:finite-smallball}. Fourier integrability justifies the
exchange of expectation and integration. The argument works conditionally
on any sigma-field supplying conditional iid variables and moment bounds;
the shift may be measurable in that sigma-field.

Condition on the first run through layer $L-1$ and the entire independent
second run. Its final feature rows $U_i\in\mathbb R^m$ are conditionally
iid. The initialized covariance calculation above, now choosing

\[
e_1=\min\{(8T_1)^{-1},
 \gamma/[4m\max(C_0^{\rm cov},s^2C_0^{\rm cov}+D_0^{\rm cov}T_1)]\},
\]

gives, except with probability at most $\tau/8$,

\[
\bar Q:=\mathbb E(U_iU_i^\top\mid\mathcal F)\succeq(\gamma/2)I,
\qquad \max_a\|U_{i,a}\|_{L^6(\cdot\mid\mathcal F)}\le X^4.
\]

Indeed the conditional last-layer mean entry error is at most
$s^2E_{L-1}+D_0^{\rm cov}\Delta_{L-1}\le\gamma/(4m)$.
Its required width is
$32s^2R_0^2e_1^{-2}\log[256L(m+1)^2/\tau]$, which is bounded by
the source gate in \eqref{cp:poly:eq:lower-width}. Linear growth
and Gaussian sixth moments prove the $L^6$ bound.

The squared-area and spectral argument in \Cref{cp:innovation}
proves, before its equal-diagonal specialization, that for
$\bar Q\succeq gI$, $T_q=\operatorname{tr}\bar Q$ and
$M_4=\mathbb E\|U\|^4$,
\[
\mathbb E\|U(U^\top y)-\bar Qy\|^2
 \ge\frac{[y^\top\bar Q^2(T_qI-\bar Q)y]^2}
 {4M_4\|\bar Qy\|^2}
 \ge\frac{(m-1)g^3(T_q-g)}{4M_4}\|y\|^2.
\]
Apply it conditionally with $g=\gamma/2$, $T_q\ge mg$ and
$M_4\le m^2X^{16}$. Some coordinate then has
\[
\operatorname{Var}[U_a(U^\top y)\mid\mathcal F]
 \ge\sigma_0^2:=Y^2\gamma^4/(256X^{16}).
\]
Choose the smallest maximizing index, which is $\mathcal F$-measurable.
The centered product has third absolute moment at most
$8m^3Y^3X^{24}$: use $\sum_b|y_b|\le mY$, the conditional $L^6$
bound, and $\mathbb E|V-\mathbb EV|^3\le8\mathbb E|V|^3$.
Thus $M_3^2/\sigma_0^6\le2^{30}m^6X^{96}\gamma^{-12}$.
Zero readout gives the exact identity

\[
\dot f_n(0,x_a)=\frac2{mn}\sum_iU_{i,a}(U_i^\top y).
\]

The second run and the conditional means are only a measurable shift.
Apply \eqref{cp:poly:eq:finite-smallball} with
$h=\tau\sigma_0/24$; its scalar failure is at most $\tau/4$ under
the second gate of \eqref{cp:poly:eq:lower-width}. Hence, allowing
the covariance failure, with probability at least $1-\tau/2$,

\[
\max_a|\dot f_n(0,x_a)-\dot{\widetilde f}_n(0,x_a)|
 \ge\frac{\tau Y\gamma^2}{192mX^8\sqrt n}.
\]

For each run use the training rectangle at failure $\tau/8$.
It contains the parameter-two ellipse on $[0,r]$, where the difference
is bounded by $32X^6Y/\lambda$. Substituting
$K=\lceil2\ell_n/\log2\rceil\le4\ell_n$ in
\Cref{cp:derivative-transfer} gives
\[
\|f_n-\widetilde f_n\|_{\mathcal X}
 \ge\frac{\tau Y\gamma^2}
 {6144m u_{\tau/8}X^{108}(1+\lambda)\sqrt n\ell_n^{5/2}}
       -\frac{768X^6Y}{\lambda n^2}.
\]
The tail is at most half the first term when
$n^{3/2}\ge2^{24}u_{\tau/8}X^{114}
m(1+\lambda)\ell_n^{5/2}/(\lambda\tau\gamma^2)$.
Since $\ell_n^{5/2}\le64n^{1/4}$, the third gate in
\eqref{cp:poly:eq:lower-width} suffices.
The union costs less than $\tau$ and the remaining coefficient exceeds
$a_\tau$. Zero initial predictions force a positive witness time;
the training endpoints coincide with the labels.
\end{proof}

\subsection{The third-order Legendre history estimate}
\label{cp:poly:legendre-section}

The Legendre model is exactly \eqref{cp:eq-legendre-reconstruction} and
\eqref{cp:eq-legendre-flow}, with its constant forward and zero backward
prefixes. Let

\begin{equation}
q=\left\lceil\lambda^{-3}n^{1/4}\ell_n^2e^{\sqrt{\ell_n}/2}\right\rceil,
\qquad P=X^{500}\chi^{10}.
\label{cp:poly:eq:legendre-order}
\end{equation}

For fixed positive $a,b$, define the numerical logarithm coefficient

\[
C(a,b)=\max\{1,e^b(a/(eb))^a\}.
\]

Maximizing $x^ae^{-bx}$ over $x\ge1$ proves
$\ell_n^a\le C(a,b)n^b$. Define the structural gate

\begin{equation}
N_{\rm Leg}=\left\lceil\max\left\{
(2e^{100}\lambda^{-3})^8, (4P^3)^4,
[8P^4C(2,1/64)]^8, d^4
\right\}\right\rceil.
\label{cp:poly:eq:legendre-gate}
\end{equation}

\begin{proposition}[Improved error without additional state]
\label{cp:poly:legendre}

On the finite training-source event, for $n\ge N_{\rm Leg}$,

\begin{equation}
\|f_{\rm Leg}-f_n\|_{\mathcal X}
 \le4YP^5\ell_n^2n^{-47/64}.
\label{cp:poly:eq:legendre-error}
\end{equation}

The exact moving and fixed inventories remain
$n(d+1)+1+2(L-1)mnq$ and $(L-1)n^2$, respectively.

\end{proposition}

\begin{proof}

Use the mobility coordinates $\theta$ defined above. For a finite horizon
$T_0$, let
\[
d_{T_0}=\int_0^{T_0}\sum_{j=2}^L\|\mathcal E_j(t)\|_Fdt,\qquad
D_{T_0}=\sup_{t\le T_0}\sum_{j=1}^{L+1}
 \|\widehat\theta_j(t)-\theta_j(t)\|,\qquad A=\tau(T_0),
\]
where $\mathcal E_j$ is the exact defect
\eqref{cp:eq-legendre-defect}. By \Cref{cp:legendre-fit},
\[
\widehat\rho(t)\le Ye^{-\lambda t/4},\qquad
\int_t^{T_0}\widehat\rho(s)ds\le4\widehat\rho(t)/\lambda,\qquad
A\le1+4Y/\lambda\le2.
\]
Bootstrap $D_{T_0}\le Y/\sqrt n$.
In backward and response subtraction use the dense state as reference,
so changed gates multiply its coordinate caps
\eqref{cp:poly:eq:real-response}. Forward and feature-pair differences
cost $X^5D_{T_0}$ and $X^9D_{T_0}$; the changed response gate costs
$YX^{116}\chi^3(1+\sqrt{\ell_n})D_{T_0}$.
Propagation $(9s)^L\le X^2$ and at most $L\le X$ terms give
\[
\begin{aligned}
\max_{a,j}\|\widehat k_a^{(j)}-k_a^{(j)}\|_2/\sqrt n
 &\le X^{40}(1+\sqrt{\ell_n})D_{T_0},\\
\max_{a,b,j}\|\widehat R_{ba}^{(j)}-R_{ba}^{(j)}\|_2/\sqrt n
 &\le X^{130}\chi^4(1+\sqrt{\ell_n})D_{T_0}.
\end{aligned}
\]
One coordinate conversion and the bootstrap therefore yield
\[
\max|\widehat k|\le YX^{50}\chi(1+\sqrt{\ell_n}),\qquad
\max|\widehat R|\le YX^{135}\chi^4(1+\sqrt{\ell_n}).
\]

Put $\widehat v_a=\widehat r_a/\widehat\rho$ and let $\mathcal F$
be the dense vector field at the reconstructed closure state.
Finite-time positivity of $\widehat\rho$ and stationary continuation
at residual zero follow from \Cref{cp:legendre-fit}. The ordinary
feature velocity per unit activity is
\[
a_b^{(j)}
=D_\theta h_b^{(j)}(\widehat\theta)\mathcal F(\widehat\theta)/\widehat\rho
=-\frac2m\phi'(\widehat z_b^{(j)})\odot
                 \sum_a\widehat v_a\widehat R_{ba}^{(j)}.
\]
The Legendre residual calculation gives
$\|\widehat v\|_2/\sqrt m=1$ and
$\|\dot{\widehat v}\|_2/\sqrt m\le X^{11}$.
Differentiating the actual closure trajectory, including its defect, gives
\begin{equation}
\begin{aligned}
\max_{b,j}\|\dot{\widehat z}_b^{(j)}\|_2/\sqrt n
 &\le YX^{30}\chi^3\widehat\rho,\\
\max_{a,j}\|\dot{\widehat\delta}_a^{(j)}\|_2/\sqrt n
 &\le X^{110}\chi^4(1+\sqrt{\ell_n})\widehat\rho,\\
\max_{a,b,j}\|\dot{\widehat R}_{ba}^{(j)}\|_2/\sqrt n
 &\le X^{220}\chi^8(1+\sqrt{\ell_n})\widehat\rho,\\
\max_{b,j}(\|a_b^{(j)}\|_2+\|\dot a_b^{(j)}\|_2)/\sqrt n
 &\le YX^{230}\chi^8(1+\sqrt{\ell_n}).
\end{aligned}
\label{cp:poly:eq:legendre-derivatives}
\end{equation}
Indeed \eqref{cp:eq-legendre-pointwise-defect} bounds full hidden speed
by $YX^{20}\chi^3\widehat\rho$, and forward differentiation costs
$X^5$. In
$\dot\delta=\phi''\dot z\,k+\phi'\dot k$,
$\dot k=\dot W^\top\delta+W^\top\dot\delta$, use the carrier
coordinate cap and $\dot z$ RMS; the top speed is
$\|\dot w\|/\sqrt n\le X^4\widehat\rho$.
The forcing is at most $X^{90}\chi^4(1+\sqrt{\ell_n})\widehat\rho$,
giving the second line after propagation.
Differentiating $R=\delta c+W(\phi'R)$ gives
$\dot\delta c+\delta\dot c+\dot W(\phi'R)
 +W(\phi''\dot zR)+W(\phi'\dot R)$.
The response cap controls the curvature product; pairing derivatives
and operator bounds control the rest, with forcing below
$X^{180}\chi^8(1+\sqrt{\ell_n})\widehat\rho$.
Finally the three terms of $\dot a$ use respectively
$\dot z$ and the response cap, $\dot{\widehat v}$ and response RMS,
and $\dot R$ with $\|\widehat v\|_2/\sqrt m=1$; $\widehat\rho\le Y$
gives the last line. No coordinate $\dot z$ bound or
$\sqrt n\|\mathcal E\|$ is used.

Subtract the accumulated physical feature defect from its history:

\[
g_b^{(j)}(\tau(t))=\widehat h_b^{(j)}(t)
 -\int_0^tD_\theta h_b^{(j)}(\widehat\theta(s))\mathcal E(s)\,ds,
\]

and use its constant initial value on $[0,1]$. Aggregate history norms
use $(mn)^{-1/2}$ over samples and neurons. The forward derivative bound
gives $\|h-g\|_{L^2}/\sqrt{mn}\le X^{10}d_{T_0}$.
Exactly $\partial_\tau g=a$ and
$\partial_\tau^2g=\dot a/\widehat\rho$ on $(1,A)$.
The first derivative vanishes at the prefix join because the initial
readout vanishes. Thus no distributional jump occurs there.

For $\mathscr D_Ag=-\partial_\xi[\xi(A-\xi)\partial_\xi g]$,
the inequality $(A-\tau(t))/\widehat\rho(t)\le4/\lambda$ cancels
the apparent terminal singularity. By
\eqref{cp:poly:eq:legendre-derivatives},

\[
\|\mathscr D_Ag\|_{L^2}/\sqrt{mn}
 \le YX^{240}\chi^9(1+\sqrt{\ell_n}).
\]

The shifted Legendre modes have eigenvalues $j(j+1)$.
Two integrations by parts have zero endpoint flux and no prefix jump;
Bessel's inequality then gives
$\|(I-\Pi_q^A)g\|_{L^2}\le[q(q+1)]^{-1}\|\mathscr D_Ag\|_{L^2}$.
Projection contraction on $h-g$ yields the forward tail

\[
H_{T_0}:=\max_j\|(I-\Pi_q^A)h^{(j)}\|_{L^2}/\sqrt{mn}
 \le YP\ell_n/q^2+Pd_{T_0}.
\]

The backward-history freezing proof following
\eqref{cp:eq-legendre-ledger} uses the same closure residual and dense
response reference, hence applies unchanged. Its bounds
$b_0\le X^{42}(Y/\lambda)(1+(Y/\lambda)\sqrt{\ell_n})$,
$b_1\le X^{16}(Y/\lambda)\sqrt\chi$ and
$B_d\sqrt{4Y/\lambda}\le P\ell_n$ give, for
$b_a=(\widehat r_a/\widehat\rho)\widehat\delta_a$,
\[
B_{T_0}:=\max_j\|(I-\Pi_q^A)b^{(j)}\|_{L^2}/\sqrt{mn}
 \le YP\ell_n/q+P\ell_nD_{T_0}\le2YP\ell_n/q
 \quad(q\le\sqrt n).
\]
The growing-projection identity \eqref{cp:eq-projection-growth}
and vanishing initial errors convert squared endpoint errors into final
squared history tails. Sample/activity Cauchy--Schwarz in the exact defect
gives $d_{T_0}\le2(L-1)H_{T_0}B_{T_0}$.
The signed comparison and ledger
\cref{cp:eq-legendre-signed,cp:eq-legendre-ledger} give
\[
D_{T_0}\le\mathcal A_nd_{T_0},\qquad
\mathcal A_n\le eX
 e^{X^{36}(Y/\lambda)^2\sqrt{\ell_n}}
 \le X^2n^{1/64}\le Pn^{1/64}.
\]
Here $c=X^{36}(Y/\lambda)^2\le X^{-24}$ and
$c\sqrt{\ell_n}\le\ell_n/64+16c^2$.
For $q\ge4P^3\ell_n$, absorb the term linear in $d_{T_0}$,
using $Y\le1$ and $2(L-1)\le P$, to obtain
\begin{equation}
d_{T_0}\le4YP^3\ell_n^2/q^3,\qquad
D_{T_0}\le4YP^4\ell_n^2n^{1/64}/q^3.
\label{cp:poly:eq:third-order-defect}
\end{equation}
The gate \eqref{cp:poly:eq:legendre-gate} ensures
$q\le\sqrt n$, $q\ge4P^3\ell_n$, and
$8P^4\ell_n^2\le n^{15/64}$.
For the first use
$\ell_n^2e^{\sqrt{\ell_n}/2}\le e^{100}n^{1/8}$ and absorb the ceiling
by two. The second follows from $n^{1/4}\ge4P^3$.
For the third use $\ell_n^2\le C(2,1/64)n^{1/64}$.
Consequently the bootstrap improves to $D_{T_0}\le Y/(2\sqrt n)$.
Continuity excludes a first exit. All constants are independent of $T_0$,
so the estimate extends to all times and to the fitted limits. Real
whole-sphere prediction subtraction costs at most $X^7D_{T_0}$,
proving \eqref{cp:poly:eq:legendre-error}.
The exact inventory is \eqref{cp:eq-legendre-count}; $n\ge d^4$
absorbs the first-matrix term in the original headline count.

\end{proof}

\subsection{Taylor accuracy, storage and final width gates}
\label{cp:poly:assembly}

Use \Cref{cp:selected}'s exact-metric selected dynamics.
The finite source lemma supplies its real fitting, RMS, carrier and
approximation premises. Its explicit ledger gives
$F\le X^6$, $K_1\le X^{18}$, $H_C\le2X^3$, $F_c\le X^{15}$ and
$\mathcal B_n/(Y/\lambda)\le X^{22}(1+\sqrt{\ell_n})$.
Thus $e^v-1\le ve^v$ permits $C_{\rm sel}=X^{60}$ for both
coefficients of \eqref{cp:smallcap-comparison}.
Choose its exact tolerance
\[
\eta=\min\left\{1,Y,16Y/\lambda,
\frac{\lambda}
 {2C_{\rm sel}n(1+\lambda^{-1/2})(1+\sqrt{\ell_n})e^{\sqrt{\ell_n}}}
\right\}.
\]
The finite-time error is at most $Y/(2n)$; the fitted tail is at most
$Y/(2n)$ under the label-independent gate
\begin{equation}
n^7\ge2C_{\rm sel}\lambda^{-1}(1+\lambda^{-1/2}).
\label{cp:poly:eq:taylor-tail-gate}
\end{equation}
since $ne^{-8\ell_n}\le n^{-7}$.

Apply the source construction in the proof of \Cref{cp:panel} with
our disk cover and $M_n=X^6\sqrt n$. Its real interval lies strictly
inside the open disk union, so compactness supplies the rectangle
required by \Cref{cp:jets}. For declared-input span dimension $k$, use
\[
K=\max\{0,\lceil\log_2(8M_n/\eta)\rceil\},\qquad
R=2m+k+1+2(2m+p)J(K+1).
\]
The proof there gives Taylor tail $\eta/8$, base coefficient accuracy
$\eta/128$, summed base error $\eta/64$, and image error $\eta/8$
using exact initialized images and operator norm eight.
Its input-span reduction preserves panel trajectories and Gaussian
first-row laws without passive labels. The exact retained count is
\[
1020(L+1)R^2+36R+10m(k+1)+pk+(m+p)+3(m+p)d.
\]
The additional gate

\begin{equation}
n\ge N_{\rm count}:=
 [32X^6(1+C_{\rm sel})\lambda^{-3/2}(1+Y^{-1})]^4
\label{cp:poly:eq:taylor-count-gate}
\end{equation}

implies $K+1\le8\ell_n$. To check it, the logarithm defining $K$ is
at most $(3/2)\log n+\sqrt{\ell_n}+\log(1+\sqrt{\ell_n})$ plus
the logarithm of the bracket in this gate. The latter is at most
$(1/4)\log n$, and the source gate gives $\ell_n\ge16$, so each
intervening term is at most $\ell_n/4$. Division by $\log2$ and the
ceiling leave the stated bound. The branches $\eta=Y$ and $\eta=1$
are smaller. Squaring the rank bound with
\eqref{cp:poly:eq:panel-count} and $k\le m+p$ yields

\[
R^2\le X^{810}(m+p)^2
 \left[(Y/\lambda)^4\ell_n^3+\lambda^{-2}\ell_n^2b_n^2\right].
\]

The exact inventory is therefore bounded by the Taylor headline with,
for example, coefficient $X^{820}$. The inverse-label gate belongs to
this simplification only. Without it the exact $K$ displays the extra
$\log(1/Y)$ dependence. Dense arrays, jets, coefficients and proof-panel
anchors are discarded after initialization, exactly as before.

We finally prove \Cref{cp:poly:width}. Allocate $\tau=\delta/2$
to the dense lower event and $\rho=\delta/4$ to the relevant source event.
On their intersection the Legendre ratio is at most

\[
4P^5a_\tau^{-1}\ell_n^{9/2}n^{-15/64}
 \le4P^5a_\tau^{-1}C(9/2,7/64)n^{-1/8}.
\]

Consequently the following fully specified maximum is a sufficient
Legendre width:

\begin{equation}
\left\lceil\max\left\{
 N_{\rm src}(m,\delta/4),N_{\rm low}(\delta/2),N_{\rm Leg},
 [4P^5C(9/2,7/64)/(a_{\delta/2}\varepsilon)]^8
 \right\}\right\rceil.
\label{cp:poly:eq:legendre-final-width}
\end{equation}

For Taylor the absolute error $Y/n$ and the same lower event give ratio
at most $a_\tau^{-1}\ell_n^{5/2}/\sqrt n\le64a_\tau^{-1}n^{-1/4}$.
A sufficient width is

\begin{equation}
\left\lceil\max\left\{
N_{\rm src}(m+p,\delta/4),N_{\rm low}(\delta/2),N_{\rm count},
[2C_{\rm sel}\lambda^{-1}(1+\lambda^{-1/2})]^{1/7},
[64/(a_{\delta/2}\varepsilon)]^4
\right\}\right\rceil.
\label{cp:poly:eq:taylor-final-width}
\end{equation}

The probability union costs less than $\delta$; no independence between
the source and comparison events is required. Since
$\lambda^{-1}=m/\gamma$, $u_\rho$ is logarithmic in $mL/\rho$,
$L\le X$, and $a_\tau^{-1}$ has only the displayed polynomial factors,
these maxima are bounded by the polynomial asserted in the theorem.
The proof counts retained real coordinates and exact-arithmetic error;
it does not infer efficient preprocessing or finite precision from a
polynomial sufficient width.
